\section{Processi primari}

\subsection{Fornitura}

\subsubsection{Introduzione}
Conformemente allo standard ISO/IEC 12207:1995, il processo di fornitura definisce un insieme strutturato di attività, metodi e procedure orientate a garantire la fornitura del prodotto software richiesto dal proponente Halue S.r.l. Società Benefit. \\
In particolare, il gruppo \textbf{Git Happens} focalizza questo processo sul monitoraggio e sul coordinamento delle attività necessarie per la realizzazione del progetto ``Sonar'', assicurando che il prodotto finale soddisfi i requisiti specificati e venga consegnato nei tempi e nei costi previsti.
\\ \\
Il processo prende avvio con l'analisi approfondita del capitolato d'appalto e dei vincoli tecnologici imposti dall'azienda proponente. Successivamente, si procede alla formalizzazione degli accordi e alla pianificazione dettagliata delle risorse e delle attività tramite il documento Piano di Progetto.

\subsubsection{Documentazione fornita}
\begin{itemize}
    \item \textbf{Valutazione dei capitolati}: documento Valutazione Capitolati raccoglie l’analisi delle diverse proposte, con lo scopo di selezionare la soluzione più in linea con le necessità del gruppo;
    \item \textbf{\term{Analisi dei Requisiti}}: documento contenente la descrizione dettagliata dei casi d'uso, dei requisiti funzionali e di qualità del sistema Sonar;
    \item \textbf{Piano di Progetto}: documento di pianificazione che include l'analisi dei rischi, la suddivisione dei ruoli, il preventivo e il consuntivo di periodo;
    \item \textbf{Piano di Qualifica}: documento che illustra le strategie di verifica, validazione e le metriche di qualità adottate dal gruppo;
    \item \textbf{Specifica Tecnica}: documento orientato all'architettura di riferimento (basata su \term{Next.js}, \term{Vercel Workflows} e \term{Supabase}) e alle scelte di design adottate;
    \item \textbf{Glossario}: raccolta dei termini tecnici e degli acronimi per garantire un'interpretazione univoca della documentazione;
    \item \textbf{Lettera di Presentazione}: documento di accompagnamento per le revisioni di progetto (\term{RTB}).
\end{itemize}

\subsection{\term{Sviluppo}}

\subsubsection{Introduzione}
Il processo di sviluppo, in accordo con lo standard ISO/IEC 12207:1995, comprende l'insieme delle attività di analisi, progettazione, codifica, testing e configurazione dell'ambiente. \\
Il gruppo \textbf{Git Happens} adotta un approccio incrementale e iterativo, garantendo la tracciabilità tra i requisiti e l'architettura implementata per il sistema Sonar.
\subsection{Analisi dei requisiti}

\subsubsection{Descrizione}
L'attività di analisi dei requisiti costituisce un momento fondamentale nello sviluppo del progetto Sonar, commissionato dall'azienda Halue S.r.l. SB. Essa ha lo scopo di definire in modo rigoroso, completo e privo di ambiguità le esigenze del proponente, stabilendo le basi per la successiva progettazione, codifica e verifica del sistema. L'analisi si concentra sulla comprensione approfondita del dominio applicativo degli store Shopify enterprise e sulla formalizzazione dei casi d'uso, delle funzionalità e dei vincoli operativi.

\subsubsection{Obiettivi}
Gli obiettivi principali dell'attività di analisi dei requisiti comprendono:
\begin{itemize}
    \item Definire con precisione gli obiettivi e le funzionalità del prodotto Sonar in stretta condivisione con l'azienda proponente, rispecchiandone le aspettative;
    \item Identificare, documentare e validare i requisiti funzionali, qualitativi e di vincolo del sistema;
    \item Facilitare la comprensione comune tra tutti gli stakeholder del progetto (membri del team, docenti e proponente);
    \item Permettere una stima accurata delle tempistiche e delle risorse necessarie;
    \item Fornire ai progettisti specifiche chiare e strutturate per orientare l'architettura software;
    \item Favorire l'attività di verifica e testing definendo riferimenti pratici e verificabili.
\end{itemize}

\subsubsection{Documentazione}
Il gruppo redige il documento \emph{Analisi dei Requisiti}, il quale costituisce il riferimento ufficiale per le specifiche del sistema. Esso contiene:
\begin{itemize}
    \item \textbf{Introduzione e Obiettivi:} descrizione dello scopo del documento e panoramica degli obiettivi del prodotto Sonar;
    \item \textbf{Descrizione del prodotto:} analisi approfondita del contesto, degli obiettivi generali, delle funzionalità core (acquisizione, normalizzazione a grafo, analisi statica del tema, motore di rilevamento orfani e presentazione) e dei vincoli di dominio (operatività in sola lettura, esclusione del monitoraggio real-time e dei dati di business);
    \item \textbf{Caratteristiche degli utenti:} profilazione delle figure professionali di riferimento (sviluppatori Shopify, architetti software, project manager tecnici);
    \item \textbf{Tecnologie di riferimento:} quadro tecnologico basato su Next.js, ecosistema Vercel (Workflows, Queues, Sandbox, Blob) e Supabase con Row Level Security (RLS);
    \item \textbf{Casi d'uso:} modellazione dettagliata degli scenari d'interazione tra gli attori e il sistema;
    \item \textbf{Requisiti:} classificazione puntuale dei requisiti in obbligatori (RFO), desiderabili (RFD) e opzionali (RFOp).
\end{itemize}

\subsubsection{Casi d'uso}
I casi d'uso forniscono una descrizione dettagliata delle funzionalità del sistema dal punto di vista degli utenti e delle entità esterne, illustrando come la piattaforma risponde a determinate azioni o scenari operativi. 
Ogni caso d'uso è identificato da un codice univoco secondo la sintassi \texttt{UC[Numero caso d'uso].[Numero sottocaso d'uso] - [Titolo]} ed è corredato da una scheda testuale strutturata nei seguenti campi:
\begin{enumerate}
    \item \textbf{Identificativo e Titolo:} codice univoco e descrizione sintetica dell'obiettivo funzionale;
    \item \textbf{Attore principale:} entità esterna che avvia e guida attivamente l'interazione con il sistema;
    \item \textbf{Descrizione:} sintesi esplicativa della funzionalità modellata;
    \item \textbf{Precondizioni:} stato in cui deve trovarsi il sistema affinché l'interazione possa avere inizio;
    \item \textbf{Postcondizioni:} stato risultante del sistema al termine della corretta esecuzione;
    \item \textbf{Scenario principale:} sequenza ordinata e numerata delle azioni e risposte;
    \item \textbf{Estensioni:} eventuali flussi alternativi o di gestione degli errori;
    \item \textbf{User story associata:} formulazione sintetica della funzionalità dal punto di vista dell'utente nella forma \emph{``Come [utente] desidero poter [funzionalità] per [valore aggiunto]''}.
\end{enumerate}
Gli attori identificati nel sistema comprendono l'utente non autenticato, l'utente autenticato (sviluppatore o tech lead associato a un tenant) e il temporizzatore (scheduler) per l'avvio delle scansioni periodiche.

\subsubsection{Diagrammi dei casi d'uso}
I diagrammi dei casi d'uso sono strumenti di modellazione grafica che permettono di visualizzare in modo chiaro, intuitivo e immediato le funzionalità offerte dal sistema e le relative interazioni con gli attori esterni, senza approfondire i dettagli implementativi di basso livello.

\subsubsection{Requisiti}
La sezione dei requisiti definisce in modo strutturato l'insieme delle richieste, delle specifiche e dei vincoli operativi che il software deve soddisfare. Essi si suddividono in base alla loro priorità e natura:
\begin{itemize}
    \item \textbf{Requisiti funzionali obbligatori:} rappresentano il nucleo essenziale e irrinunciabile del sistema, costituendo la base fondamentale su cui viene valutata la realizzazione del progetto.
    \item \textbf{Requisiti funzionali desiderabili:} comprendono funzionalità non strettamente indispensabili per il funzionamento primario, ma che offrono un significativo valore aggiunto all'utente.
    \item \textbf{Requisiti funzionali opzionali:} raggruppano caratteristiche secondarie o migliorative la cui implementazione può essere valutata o eventualmente contrattata in base all'effettivo avanzamento temporale e delle risorse.
\end{itemize}

\subsubsection{Strumenti}
PlantUML (tramite estensione dedicata per Visual Studio Code): strumento utilizzato dal team per la generazione e la gestione visuale dei diagrammi dei casi d'uso in modalità \emph{Text-to-UML}.

\subsection{Progettazione}
La fase di progettazione definisce l'architettura logica e di deployment del sistema Sonar, separando nettamente i moduli di acquisizione, normalizzazione, analisi e presentazione. \\
\\
I principali aspetti affrontati comprendono:
\begin{itemize}
    \item \textbf{Modello a Grafo}: strutturazione delle entità Shopify e delle relative dipendenze all'interno di un database relazionale Supabase (PostgreSQL) protetto da \term{Row Level Security} (RLS);
    \item \textbf{\term{Diagrammi UML} e C4}: utilizzo di diagrammi architetturali e di classe per descrivere la componentistica del sistema;
    \item \textbf{Design Pattern}: adozione di soluzioni architetturali orientate all'estensibilità del motore di regole per il rilevamento di codice morto e componenti orfani;
    \item \textbf{Test}: pianificazione dei test di unità, integrazione e sistema tramite framework dedicati come Vitest e Playwright.
\end{itemize}

\subsection{Codifica}
L'attività di codifica è affidata ai programmatori del gruppo \textbf{Git Happens}, i quali traducono le specifiche tecniche in codice sorgente robusto e manutenibile scritto in \term{TypeScript} (strict mode). Vengono rigorosamente seguite le norme di codifica ispirate ai principi di "Clean Code", garantendo la pulizia del codice, la modularità e l'assenza di errori di linting in fase di integrazione continua (\term{CI}). \\
Come strumenti di sviluppo vengono impiegati \term{Visual Studio Code} e la piattaforma \term{GitHub} per il versionamento.

\subsection{Configurazione ambiente}
La configurazione dell'ambiente di esecuzione e di sviluppo del progetto Sonar si basa sull'ecosistema Vercel e Supabase. Nello specifico, il gruppo gestisce:
\begin{itemize}
    \item \textbf{Vercel Workflows e Queues}: per l'orchestrazione durevole e la ripresa delle scansioni in seguito a limiti di rate limiting o interruzioni;
    \item \textbf{Containerizzazione e servizi cloud}: utilizzo di configurazioni standard per l'esecuzione isolata tramite Vercel Sandbox e l'archiviazione di snapshot su Vercel Blob;
    \item \textbf{Automazione}: integrazione di flussi \term{CI}/CD tramite \term{GitHub} Actions per l'esecuzione automatica dei test e la verifica della copertura del codice.
\end{itemize}